% Generated draft: requires mathematical and source review.
\documentclass[11pt]{article}
\usepackage[margin=25mm]{geometry}
\usepackage{fontspec}
\setmainfont{Latin Modern Roman}
\usepackage{amsmath,amssymb,mathtools,unicode-math}
\setmathfont{Latin Modern Math}
\usepackage{enumitem,needspace,xurl,hyperref}
\setlength{\parindent}{0pt}
\setlength{\parskip}{5pt}
\setlength{\emergencystretch}{3em}
\newcommand{\N}{\mathbb{N}}
\newcommand{\Z}{\mathbb{Z}}
\newcommand{\Q}{\mathbb{Q}}
\newcommand{\R}{\mathbb{R}}
\newcommand{\C}{\mathbb{C}}
\newcommand{\F}{\mathbb{F}}
\newcommand{\A}{\mathbb{A}}
\newcommand{\PP}{\mathbb{P}}
\newcommand{\E}{\mathbb{E}}
\newcommand{\eps}{\varepsilon}
\newcommand{\dd}{\,\mathrm{d}}
\begin{document}
\section*{UnsolvedMath ID: 30005784}
\textbf{Source problem code:} catalogue entry
\subsection{Definitions and mathematical statement}
Let \(f\colon X\to Y\) be a threefold flopping contraction in the intended setup of the supplied source. Write \(\mathcal D\) for its associated triangulated category and
\[
B=\widehat E\setminus\widehat{\mathcal H}
\]
for the associated complexified hyperplane complement, where \(\widehat{\mathcal H}\) denotes the removed hyperplane locus. The supplied record does not specify the construction of \(\mathcal D\), \(\widehat E\), or \(\widehat{\mathcal H}\), or the full hypotheses on \(f\). Thus the questions below are relative to the intended source setup, rather than a fully specified class of contractions.

Let \(\operatorname{Stab}(\mathcal D)\) denote the space of Bridgeland stability conditions on \(\mathcal D\), and let
\[
S=\operatorname{Stab}^{\circ}(\mathcal D)
\]
be the distinguished connected component referred to in the source. The precise convention for stability conditions and any restriction on their central charges are not supplied. For a base point \(b\in B\), the supplied excerpt attributes to Theorem 2 a mutation homomorphism
\[
\phi\colon\pi_1(B,b)\longrightarrow\operatorname{Aut}(\mathcal D)
\]
and a regular covering map
\[
p\colon S\longrightarrow B,
\qquad \operatorname{Deck}(p)\cong\operatorname{im}(\phi).
\]
Here \(\operatorname{Aut}(\mathcal D)\) denotes the intended group of categorical autoequivalences; its precise conventions are not specified. The excerpt states that monodromy around a root hyperplane corresponds to the mutation autoequivalences of its reference [13]. It further states that, for a threefold flopping contraction, these mutations have a geometric interpretation through Bridgeland--Chen flop functors. These assertions are supplied background, not independently verified here.

\textbf{Questions.} For the intended threefold-flop setup, determine the answers to the following questions:
\begin{enumerate}
\item Is \(p\colon S\to B\) the universal covering map?
\item Is the mutation action faithful, that is,
\[
\ker(\phi)=\{1\}?
\]
The source presents this as equivalent to the preceding question. That equivalence depends on the intended relationship between \(\phi\) and covering monodromy; an abstract identification of the deck group with \(\operatorname{im}(\phi)\) alone does not specify that relationship.
\item Is the entire stability manifold connected, equivalently,
\[
\operatorname{Stab}(\mathcal D)=S?
\]
\item Is \(S\) contractible, meaning homotopy equivalent to a point?
\end{enumerate}

The reconstructed statement also asks whether \(B\) is a \(K(\pi,1)\) space, with \(\pi=\pi_1(B,b)\). In the usual setting of connected spaces admitting universal covers and having CW homotopy type, this means that the universal cover is contractible, equivalently that \(\pi_n(B,b)=0\) for every integer \(n\geq2\). If \(p\) is the universal cover, this is equivalent to contractibility of \(S\). Without that qualification, the record's wording ``equivalently'' does not justify identifying the \(K(\pi,1)\) question with contractibility of this particular cover, or with the entire list of questions, especially global connectedness of \(\operatorname{Stab}(\mathcal D)\).

The original excerpt also gives a parallel formulation for a category \(\mathcal C\), with complement \(E\setminus\mathcal H\), a homomorphism to \(\operatorname{Aut}(\mathcal C)\), and a regularly covering component \(\operatorname{Stab}^{\circ}(\mathcal C)\). Its Remark 3 mentions spherical twist functors, citing [9], when \(f\) is the minimal resolution of a Kleinian singularity. The constructions of these parallel objects are likewise absent. The reconstructed target supplied for this record is the threefold-flop category \(\mathcal D\). The original third question is truncated, so its full wording cannot be recovered from that excerpt alone.
\subsection{Short English statement}
For the category associated with a threefold flop, ask whether the distinguished component of its Bridgeland stability space is the universal cover of the associated hyperplane complement, whether mutations act faithfully, whether all stability conditions lie in that component, and whether the component is contractible. The supplied record omits essential definitions and geometric hypotheses. Its \(K(\pi,1)\) reformulation requires qualifications, including identifying the distinguished cover as universal.
\subsection{Sources}
\begin{itemize}
\item Supplied local UnsolvedMath registry record, permanent ID \(30005784\), title ``Stability Manifolds for Threefold Flops,'' exact source problem code \texttt{catalogue entry}, in the Oberwolfach Reports Open Problems collection. No individual proposer is identified; the supplied proposed year is 2024.
\item Supplied source citation: \emph{Mini-Worskhop: Artin Groups meet Triangulated Categories} (2024), Oberwolfach Reports, pp.~203--234, DOI \texttt{10.4171/owr/2024/4}. Source URL and supplied literature-source link labelled ``original OWR report'': \url{https://doi.org/10.4171/owr/2024/4}. The record describes these reports as documenting mathematical workshops and open-problem discussions. The supplied extract refers to Theorem 2 and Remark 3, and to internal references [13] and [9]; bibliographic details for those internal references are not supplied.
\item Supplied literature assessment, checked \texttt{2026-08-22}: status labelled ``open.'' The assessment reports that a connected component covering the hyperplane complement is proved, while faithfulness, universal-cover status, global connectedness, and contractibility are posed as unresolved questions. This is supplied dated curation derived from the OWR record, not independently verified; time-sensitive status and cited claims require independent verification.
\item The registry labels the statement \texttt{corrected\_verified} and says it was reconstructed and checked against the official Oberwolfach report. These are supplied provenance claims, not verification performed for this draft.
\end{itemize}
\end{document}
